\chapter{Logika, Himpunan dan Pemetaan}\label{ch:b1:logic}

Sampai sekarang, pembuktian dijalankan dengan gagasan yang informal
namun jujur tentang apa arti ``membuktikan''. Bab pertama matematika
tingkat sarjana ini menyatakan aturan mainnya secara tegas: apa itu
\omterm{def:b1:logic:statement}{pernyataan} matematis, bagaimana perangkai dan kuantor menggabungkan
\omterm{def:b1:logic:statement}{pernyataan}, langkah mana yang sah di dalam sebuah bukti --- lalu, di
atas landasan itu, membangun dua bahasa universal matematika: \omterm{def:b1:logic:sets}{himpunan}
dan \omterm{def:b1:logic:map}{pemetaan}.

\section{Pernyataan dan perangkai}

\begin{definition}[Pernyataan, perangkai]\label{def:b1:logic:statement}
Sebuah \emph{pernyataan}\index{pernyataan} (atau \emph{proposisi})
adalah kalimat yang bernilai benar (B) atau salah (S) --- tepat satu di
antara keduanya. Dari pernyataan $P$ dan $Q$ dibentuk:
\begin{itemize}
  \item \emph{negasi} $\lnot P$ (``bukan $P$''), benar tepat ketika
        $P$ salah;
  \item \emph{konjungsi} $P \land Q$ (``$P$ dan $Q$''), benar tepat
        ketika keduanya benar;
  \item \emph{disjungsi} $P \lor Q$ (``$P$ atau $Q$''), benar tepat
        ketika sekurang-kurangnya satu benar (``atau'' di sini bersifat
        inklusif);
  \item \emph{implikasi} $P \implies Q$, salah tepat ketika $P$
        benar dan $Q$ salah;
  \item \emph{ekuivalensi} $P \iff Q$, benar tepat ketika $P$ dan
        $Q$ mempunyai nilai kebenaran yang sama.
\end{itemize}
\end{definition}

\begin{remark}
Tabel kebenaran $P \implies Q$ layak dihentikan sejenak: ketika $P$
salah, $P \implies Q$ bernilai \emph{benar}, apa pun $Q$-nya. ``Jika
$2 < 1$ maka $0 = 5$'' adalah implikasi yang benar. Sebuah implikasi
tidak mengatakan apa-apa tentang keadaan ketika hipotesisnya gagal.
\end{remark}

\begin{proposition}[Kaidah perhitungan pernyataan]\label{prop:b1:logic:rules}
Untuk setiap \omterm{def:b1:logic:statement}{pernyataan} $P$, $Q$, $R$:
\begin{enumerate}
  \item $\lnot(\lnot P) \iff P$;
  \item hukum De Morgan\index{hukum De Morgan}:
        $\lnot(P \land Q) \iff (\lnot P) \lor (\lnot Q)$ dan
        $\lnot(P \lor Q) \iff (\lnot P) \land (\lnot Q)$;
  \item $(P \implies Q) \iff \bigl((\lnot P) \lor Q\bigr)$, sehingga
        $\lnot(P \implies Q) \iff P \land (\lnot Q)$;
  \item kontraposisi\index{kontraposisi}:
        $(P \implies Q) \iff \bigl((\lnot Q) \implies (\lnot P)\bigr)$;
  \item $(P \iff Q) \iff \bigl((P \implies Q) \land (Q \implies
        P)\bigr)$;
  \item sifat distributif: $P \land (Q \lor R) \iff (P \land Q) \lor
        (P \land R)$ dan $P \lor (Q \land R) \iff (P \lor Q) \land
        (P \lor R)$.
\end{enumerate}
\end{proposition}

\begin{proof}
Setiap ekuivalensi diperiksa dengan membandingkan tabel kebenaran: dua
\omterm{def:b1:logic:statement}{pernyataan} majemuk yang dibangun dari $P$, $Q$, $R$ ekuivalen tepat
ketika keduanya bernilai kebenaran sama pada masing-masing (empat atau
delapan) kasus. Kita tuliskan satu tabel selengkapnya, untuk hukum De
Morgan yang pertama:
\begin{center}
\begin{tabular}{cc|c|c|cc|c}
$P$ & $Q$ & $P \land Q$ & $\lnot(P \land Q)$ & $\lnot P$ & $\lnot Q$ &
$(\lnot P) \lor (\lnot Q)$ \\
\hline
B & B & B & S & S & S & S \\
B & S & S & B & S & B & B \\
S & B & S & B & B & S & B \\
S & S & S & B & B & B & B \\
\end{tabular}
\end{center}
Kolom $4$ dan $7$ berimpit, dan itulah bukti hukum tersebut. Untuk
kontraposisi ada jalan pintas lisan yang lebih cepat: $P \implies Q$
salah tepat pada kasus ($P$ benar, $Q$ salah), sedangkan
$(\lnot Q) \implies (\lnot P)$ salah tepat pada kasus ($\lnot Q$
benar, $\lnot P$ salah), yaitu ($Q$ salah, $P$ benar) --- satu kasus
yang persis sama, sehingga kedua implikasi itu bertabel sama. Kaidah
selebihnya diperiksa dengan cara yang sama; perhatikan bahwa (3)
menyusutkan setiap implikasi menjadi disjungsi, sehingga (2) secara
mekanis melahirkan kaidah negasi $\lnot(P \implies Q) \iff P
\land (\lnot Q)$: untuk membantah sebuah implikasi kita harus
menunjukkan satu kasus yang hipotesisnya berlaku dan kesimpulannya
gagal.
\end{proof}

\section{Kuantor}

\begin{definition}[Kuantor]\label{def:b1:logic:quantifiers}
Misalkan $P(x)$ sifat dari suatu unsur $x$ pada \omterm{def:b1:logic:sets}{himpunan} $E$.
\begin{itemize}
  \item $\forall x \in E,\ P(x)$ (``untuk setiap $x$ di $E$, $P(x)$'')
        benar bila setiap unsur $E$ memenuhi $P$;
  \item $\exists x \in E,\ P(x)$ (``terdapat $x$ di $E$ sedemikian
        sehingga $P(x)$'') benar bila sekurang-kurangnya satu unsur
        $E$ memenuhi $P$.
\end{itemize}
Kita tulis $\exists!$ untuk ``terdapat dengan tunggal''.
\end{definition}

\begin{proposition}[Negasi kuantor]\label{prop:b1:logic:negquant}
\[
\lnot\bigl(\forall x \in E,\ P(x)\bigr) \iff
\exists x \in E,\ \lnot P(x),
\qquad
\lnot\bigl(\exists x \in E,\ P(x)\bigr) \iff
\forall x \in E,\ \lnot P(x).
\]
\end{proposition}

\begin{proof}
Kita bahas ekuivalensi pertama dalam dua arah; yang kedua bersifat
simetris. Jika $\forall x \in E,\ P(x)$ salah, maka tidak setiap unsur
memenuhi $P$: \omterm{def:b1:logic:sets}{himpunan} $A = \{x \in E : \lnot P(x)\}$ tidak mungkin
kosong, dan unsur mana pun darinya menjadi saksi bagi $\exists x \in
E,\ \lnot P(x)$. Sebaliknya, jika ada $x_0 \in E$ yang memenuhi $\lnot
P(x_0)$, maka $x_0$ adalah contoh penyangkal dan \omterm{def:b1:logic:statement}{pernyataan} universal
itu gugur. Untuk kaidah kedua: ``tidak ada $x$ yang memenuhi $P$''
berarti \omterm{def:b1:logic:sets}{himpunan} $\{x : P(x)\}$ kosong, yaitu setiap $x$ terletak di
komplemennya $A$. Diterapkan berantai pada rangkaian kuantor
bersarang, kedua kaidah itu memberikan prosedur mekanis
\cref{ex:b1:logic:limit}: negasi berjalan dari kiri ke kanan,
membalik setiap $\forall$ menjadi
$\exists$ dan setiap $\exists$ menjadi $\forall$, lalu akhirnya
menegasikan predikat yang terdalam.
\end{proof}

\begin{example}[Menegasikan kalimat matematis sehari-hari]\label{ex:b1:logic:negpractice}
Misalkan $f \colon \R \to \R$. Kalimat ``$f$ naik'' berbunyi
\[
\forall x \in \R,\ \forall y \in \R,\quad
x \leq y \implies f(x) \leq f(y) ,
\]
dan negasinya, menurut \cref{prop:b1:logic:negquant} beserta kaidah
$\lnot(P \implies Q) \iff P \land \lnot Q$:
\[
\exists x \in \R,\ \exists y \in \R,\quad
x \leq y \ \text{ dan }\ f(x) > f(y) :
\]
satu pasang saksi sudah cukup. Demikian pula ``$f$ terbatas'' berarti
$\exists M \in \R,\ \forall x \in \R,\ \abs{f(x)} \leq M$, dengan
negasi
\[
\forall M \in \R,\ \exists x \in \R,\quad \abs{f(x)} > M :
\]
\emph{berapa pun} batas yang diajukan, ada titik yang melampauinya.
Inti gagasannya: negasi yang benar tidak pernah memuat ``bukan'' yang
dikenakan pada blok berkuantor --- ia adalah \omterm{def:b1:logic:statement}{pernyataan} positif yang
baru, dengan peran yang bertukar: kini kita yang menghasilkan saksi,
padahal tadinya kita yang menerimanya.
\end{example}

\begin{example}[Urutan kuantor]\label{ex:b1:logic:order}
\omterm{def:b1:logic:order}{Urutan} kuantor yang berlainan jenis itu penting:
\[
\forall x \in \R,\ \exists y \in \R,\ y > x
\quad\text{benar (ambil } y = x+1\text{),}
\]
\[
\exists y \in \R,\ \forall x \in \R,\ y > x
\quad\text{salah (tak ada bilangan real yang melampaui semua bilangan real).}
\]
Pada \omterm{def:b1:logic:statement}{pernyataan} pertama $y$ boleh bergantung pada $x$; pada yang
kedua, satu $y$ harus berlaku untuk semua $x$. Sebaliknya, dua kuantor
yang sejenis selalu dapat dipertukarkan.
\end{example}

\begin{example}[Membaca definisi dengan tiga kuantor]\label{ex:b1:logic:limit}
Kalimat ``barisan $(u_n)$ konvergen ke $\ell$'' akan dituliskan pada
\cref{ch:b1:seq} sebagai
\[
\forall \varepsilon > 0,\ \exists N \in \N,\ \forall n \geq N,\quad
\abs{u_n - \ell} \leq \varepsilon .
\]
Negasinya, dengan \cref{prop:b1:logic:negquant} yang diterapkan tiga kali, adalah
\[
\exists \varepsilon > 0,\ \forall N \in \N,\ \exists n \geq N,\quad
\abs{u_n - \ell} > \varepsilon .
\]
Kemampuan menegasikan kalimat semacam itu secara mekanis, tanpa
memikirkan artinya, adalah keterampilan yang sesungguhnya: ia
memisahkan kerja logika dari kerja matematika.
\end{example}

\section{Teknik pembuktian}

\begin{method}[Pola pembuktian yang baku]\label{met:b1:logic:proofs}
Untuk membuktikan\dots
\begin{enumerate}
  \item \emph{implikasi $P \implies Q$ secara langsung}: andaikan $P$,
        simpulkan $Q$;
  \item \emph{dengan kontraposisi}: andaikan $\lnot Q$, simpulkan
        $\lnot P$ --- sah menurut \cref{prop:b1:logic:rules}~(4);
  \item \emph{dengan kontradiksi}: andaikan \omterm{def:b1:logic:statement}{pernyataan} itu salah, lalu
        turunkan sebuah kontradiksi;
  \item \emph{sebuah ekuivalensi}: buktikan kedua implikasinya secara
        terpisah (atau rangkaikan ekuivalensi yang sudah dikenal);
  \item \emph{\omterm{def:b1:logic:statement}{pernyataan} ``untuk setiap''}: ambil $x$ yang
        \emph{sembarang} di $E$ (``misalkan $x \in E$'') lalu buktikan
        $P(x)$;
  \item \emph{\omterm{def:b1:logic:statement}{pernyataan} ``terdapat''}: tunjukkan sebuah saksi, atau
        buktikan keberadaannya secara tak langsung;
  \item \emph{dengan induksi}: lihat \cref{thm:b1:logic:induction}.
\end{enumerate}
Ketika membuktikan \omterm{def:b1:logic:statement}{pernyataan} tentang sebuah unsur yang dipilih dengan
cermat namun tetap sembarang, jangan sekali-kali memberi unsur itu
sifat tambahan: ``misalkan $x \in \R$'' yang disusul ``karena
$x > 0$\dots'' tidak membuktikan apa pun untuk $x$ yang negatif.
\end{method}

\begin{remark}[Jebakan yang lazim dalam pembuktian]\label{rem:b1:logic:pitfalls}
Empat jebakan klasik, semuanya pantas disebut sekali.
\begin{enumerate}
  \item \emph{Konvers, bukan kontraposisi.} $Q \implies P$
        \emph{tidak} ekuivalen dengan $P \implies Q$; hanya $\lnot Q
        \implies \lnot P$ yang ekuivalen. ``Jika hujan turun, jalan
        menjadi basah'' tidak memberi hak untuk menyimpulkan hujan
        dari jalan yang basah.
  \item \emph{Membuktikan ekuivalensi lewat satu implikasi saja.}
        Klaim ``jika dan hanya jika'' adalah dua teorema; sebutkan
        arah mana yang sedang dibuktikan, lalu buktikan keduanya.
        Rangkaian $\iff$ hanya sah bila \emph{setiap} mata rantainya
        benar-benar dapat dibalik --- mengkuadratkan persamaan,
        misalnya, tidak.
  \item \emph{Bukti yang berjalan mundur.} Berangkat dari kesimpulan
        yang diinginkan lalu menurunkan \omterm{def:b1:logic:statement}{pernyataan} yang benar tidak
        membuktikan apa-apa (dari $-1 = 1$ orang menurunkan $1 = 1$
        yang benar dengan mengkuadratkan). Sebuah perhitungan boleh
        \emph{ditemukan} secara mundur, tetapi ia harus
        \emph{dituliskan} maju, atau dengan ekuivalensi yang eksplisit.
  \item \emph{Saksi tertentu lawan unsur sembarang.} Untuk membuktikan
        $\exists x,\ P(x)$, kita boleh menunjukkan satu $x$ yang
        dipilih dengan cerdik; untuk membuktikan $\forall x,\ P(x)$,
        unsur $x$ yang dipilih harus tetap sembarang. Mencampur
        keduanya --- memeriksa klaim universal pada sebuah contoh ---
        adalah kesalahan yang paling sering muncul dalam pekerjaan
        pemula.
\end{enumerate}
\end{remark}

\begin{example}[Kontraposisi dan kontradiksi dalam praktik]\label{ex:b1:logic:sqrt2}
\emph{Untuk $n \in \N$: jika $n^2$ genap maka $n$ genap.} Dengan
kontraposisi: jika $n$ ganjil, $n = 2k+1$, maka
$n^2 = 4k^2 + 4k + 1$ ganjil.

\emph{$\sqrt 2$ irasional.} Dengan kontradiksi: andaikan
$\sqrt 2 = p/q$ dengan $p, q \in \N^*$ dan pecahan itu sudah paling
sederhana. Maka $p^2 = 2q^2$ genap, sehingga $p$ genap (butir
sebelumnya), $p = 2r$; lalu $q^2 = 2r^2$ genap, sehingga $q$ genap ---
bertentangan dengan bentuk paling sederhana tadi.
\end{example}

\begin{theorem}[Induksi matematika]\label{thm:b1:logic:induction}
Misalkan $P(n)$ sifat dari bilangan bulat $n$. Jika
\begin{enumerate}
  \item $P(0)$ benar, dan
  \item untuk setiap $n \in \N$ berlaku $P(n) \implies P(n+1)$,
\end{enumerate}
maka $P(n)$ benar untuk setiap $n \in \N$.

\emph{Induksi kuat:} kesimpulannya tidak berubah bila (2) diganti
oleh: untuk setiap $n$, $\bigl(P(0) \land \dots \land P(n)\bigr)
\implies P(n+1)$.
\end{theorem}

\begin{proof}
Ini sifat $\N$ itu sendiri, yang setara dengan: \emph{setiap \omterm{def:b1:logic:sets}{himpunan}
bagian tak kosong dari $\N$ mempunyai unsur terkecil} (kita terima
sebagai hal yang sudah diketahui). Memang, andaikan (1) dan (2)
berlaku, lalu tulis $A = \{n \in \N : P(n)
\text{ salah}\}$. Jika $A \neq \emptyset$, ia mempunyai unsur terkecil
$m$; $m \neq 0$ menurut (1); lalu $m - 1 \notin A$, sehingga $P(m-1)$
berlaku, dan (2) memberikan $P(m)$ --- kontradiksi. Jadi
$A = \emptyset$. Untuk induksi kuat, terapkan argumen yang sama:
$P(0), \dots, P(m-1)$ semuanya berlaku karena $m$ terkecil di $A$.
\end{proof}

\begin{example}[Membuktikan keberadaan yang tunggal]\label{ex:b1:logic:uniqueexistence}
\omterm{def:b1:logic:statement}{Pernyataan} $\exists!\,x,\ P(x)$ sesungguhnya \emph{dua} \omterm{def:b1:logic:statement}{pernyataan},
yang dibuktikan terpisah: keberadaan (tunjukkan atau bangun suatu
$x_0$ dengan $P(x_0)$) dan ketunggalan (andaikan $P(x)$ dan $P(x')$,
simpulkan $x =
x'$). Contoh: \emph{ada tepat satu bilangan real $x$ dengan $x^3 + x =
2$.} Keberadaan: $x_0 = 1$ memenuhi, karena $1 + 1 = 2$. Ketunggalan:
jika $x^3 + x = x'^3 + x'$, maka
\[
0 = (x^3 - x'^3) + (x - x')
= (x - x')\,\bigl(x^2 + xx' + x'^2 + 1\bigr),
\]
dan faktor keduanya bernilai positif (ia sama dengan $\bigl(x +
\tfrac{x'}2\bigr)^2 + \tfrac34 x'^2 + 1 \geq 1$), sehingga $x = x'$.
Perhatikan pembagian tugasnya: keberadaan memakai tebakan yang
beruntung, ketunggalan memakai aljabar yang berlaku untuk penyelesaian
\emph{sembarang} --- tak satu pun di antara kedua argumen itu
mengerjakan tugas yang lain, dan melupakan paruh kedua adalah godaan
yang selalu hadir begitu sebuah penyelesaian ditemukan.
\end{example}

\begin{example}\label{ex:b1:logic:induction}
Untuk setiap $n \in \N^*$: $\;\sum_{k=1}^n k = \frac{n(n+1)}{2}$.
Basis $n = 1$: kedua ruas sama dengan $1$. Langkah: dengan
mengandaikan rumus itu berlaku untuk $n$,
\[
\sum_{k=1}^{n+1} k = \frac{n(n+1)}{2} + (n+1)
= (n+1)\Bigl(\frac n2 + 1\Bigr) = \frac{(n+1)(n+2)}{2}. \qedhere
\]
\end{example}

\begin{example}[Induksi kuat dalam praktik]\label{ex:b1:logic:strongind}
\emph{Setiap bilangan bulat $n \geq 2$ merupakan hasil kali bilangan
prima} (bilangan prima adalah bilangan bulat $\geq 2$ yang pembaginya
yang $\geq 1$ hanyalah $1$ dan dirinya sendiri; bilangan prima
dipelajari tersendiri pada \cref{ch:b1:arith}). Induksi biasa tak
berdaya di sini: mengetahui bahwa $95 = 5 \times 19$ terfaktorkan
tidak mengatakan apa pun tentang $96$. Induksi kuat justru pas.
Basisnya: $2$ prima, jadi ia hasil kali bilangan prima dengan satu
faktor. Langkah: ambil $n \geq 2$ dan andaikan setiap bilangan bulat
$m$ dengan $2 \leq m \leq n$ merupakan hasil kali bilangan prima. Jika
$n + 1$ prima, selesai. Jika tidak, $n + 1 = ab$ dengan
$2 \leq a, b \leq n$; menurut hipotesis kuat, $a$ dan $b$ keduanya
hasil kali bilangan prima, sehingga demikian pula $n + 1$. Inti
gagasannya: induksi kuat adalah alat yang tepat setiap kali
``alasan'' bagi $P(n+1)$ berada pada suatu peringkat sebelumnya yang
tak terduga, bukan pada peringkat $n$.
\end{example}

\section{Himpunan}

\begin{definition}[Operasi himpunan]\label{def:b1:logic:sets}
Gagasan \emph{himpunan}\index{himpunan} dan relasi keanggotaan
$x \in E$ kita terima sebagai gagasan primitif. Untuk himpunan $A, B$
di dalam himpunan semesta $E$:
\begin{itemize}
  \item \emph{inklusi}: $A \subseteq B$ bila $\forall x,\ x \in A
        \implies x \in B$; kesamaan $A = B$ bila $A \subseteq B$ dan
        $B \subseteq A$;
  \item \emph{gabungan} $A \cup B$, \emph{irisan} $A \cap B$,
        \emph{selisih} $A \setminus B = \{x \in A : x \notin B\}$,
        \emph{komplemen} $\overline{A} = E \setminus A$;
  \item \emph{himpunan kosong} $\emptyset$, yang termuat dalam setiap
        himpunan;
  \item \emph{himpunan kuasa}\index{himpunan kuasa} $\mathcal{P}(E)$:
        himpunan semua himpunan bagian dari $E$;
  \item \emph{hasil kali} $E \times F$: himpunan pasangan terurut
        $(x, y)$ dengan $x \in E$, $y \in F$.
\end{itemize}
\end{definition}

\begin{example}[Membiasakan diri dengan himpunan kuasa]\label{ex:b1:logic:powerset}
Untuk $E = \{a, b\}$:
\[
\mathcal P(E) = \bigl\{\, \emptyset,\ \{a\},\ \{b\},\ \{a, b\}
\,\bigr\},
\]
empat unsur --- dan perhatikan disiplin tipenya: $a \in E$ tetapi
$\{a\} \in \mathcal P(E)$; \omterm{def:b1:logic:statement}{pernyataan} $a \in \mathcal P(E)$
maupun $\{a\} \subseteq \mathcal P(E)$ keduanya salah sebagaimana
tertulis (yang kedua menuntut $a$ menjadi \emph{\omterm{def:b1:logic:sets}{himpunan} bagian} dari
$E$). Diiterasikan dari ketiadaan: $\mathcal P(\emptyset) =
\{\emptyset\}$ mempunyai satu unsur, $\mathcal P(\mathcal P(\emptyset))
= \{\emptyset, \{\emptyset\}\}$ mempunyai dua, berikutnya empat ---
\omterm{def:b1:logic:sets}{himpunan} yang anggotanya \omterm{def:b1:logic:sets}{himpunan} tetaplah \omterm{def:b1:logic:sets}{himpunan} biasa, dan
\cref{ch:b1:counting} akan membenarkan pola penggandaan itu:
$\abs{\mathcal P(E)} = 2^{\abs E}$. Menjaga agar tingkatannya
($x$, $\{x\}$, $\{\{x\}\}$) tidak tertukar adalah separuh perjuangan
dalam latihan seperti \cref{exo:b1:logic:11,exo:b1:logic:12}.
\end{example}

\begin{proposition}[Aljabar himpunan]\label{prop:b1:logic:setalgebra}
Untuk \omterm{def:b1:logic:sets}{himpunan} bagian $A, B, C$ dari $E$:
\begin{enumerate}
  \item $A \cap (B \cup C) = (A \cap B) \cup (A \cap C)$ dan
        $A \cup (B \cap C) = (A \cup B) \cap (A \cup C)$;
  \item De Morgan: $\overline{A \cup B} = \overline{A} \cap
        \overline{B}$ dan $\overline{A \cap B} = \overline{A} \cup
        \overline{B}$;
  \item $A \subseteq B \iff \overline{B} \subseteq \overline{A}$.
\end{enumerate}
\end{proposition}

\begin{proof}
Setiap kesamaan itu menerjemahkan satu kaidah
\cref{prop:b1:logic:rules} lewat kamus ($\in A$ atau bukan)
$\leftrightarrow$ (\omterm{def:b1:logic:statement}{pernyataan} benar atau
salah): misalnya $x \in \overline{A \cup B} \iff \lnot(x \in A \lor x \in
B) \iff (x \notin A) \land (x \notin B) \iff x \in \overline{A} \cap
\overline{B}$. Butir (3) adalah kontraposisi. Sebagai contoh kedua,
hukum distributif yang pertama selengkapnya:
\[
x \in A \cap (B \cup C)
\iff (x \in A) \land \bigl(x \in B \lor x \in C\bigr)
\iff \bigl(x \in A \land x \in B\bigr) \lor
\bigl(x \in A \land x \in C\bigr),
\]
menurut sifat distributif pada \cref{prop:b1:logic:rules}~(6), dan
\omterm{def:b1:logic:statement}{pernyataan} terakhir itu terbaca $x \in (A \cap B) \cup (A \cap C)$.
Setiap kesamaan \omterm{def:b1:logic:sets}{himpunan} semacam ini dapat dibuktikan lewat satu
terjemahan mekanis ini --- itulah sebabnya tak satu pun perlu dihafal.
\end{proof}

\begin{method}[Membuktikan kesamaan himpunan]\label{met:b1:logic:doubleinc}
Untuk membuktikan $A = B$, buktikan kedua inklusinya: ambil $x \in A$,
tunjukkan $x \in B$; lalu ambil $x \in B$, tunjukkan $x \in A$. Cara
lain, rangkaikan ekuivalensi $x \in A \iff \dots \iff x \in B$ bila
setiap langkahnya memang benar-benar sebuah ekuivalensi.
\end{method}

\begin{omfigure}
\begin{tikzpicture}[scale=0.9]
  \draw[thick] (-2.6,-1.9) rectangle (2.6,1.9);
  \node[below right] at (-2.6,1.9) {$E$};
  \begin{scope}
    \clip (-2.6,-1.9) rectangle (2.6,1.9);
    \fill[omDef!20] (-2.6,-1.9) rectangle (2.6,1.9);
    \fill[white] (-0.7,0) circle (1.1);
    \fill[white] (0.7,0) circle (1.1);
  \end{scope}
  \draw[omDef, thick] (-0.7,0) circle (1.1);
  \draw[omDef, thick] (0.7,0) circle (1.1);
  \node[omDef] at (-1.35,0.5) {$A$};
  \node[omDef] at (1.35,0.5) {$B$};
  \begin{scope}[xshift=6.6cm]
    \draw[thick] (-2.6,-1.9) rectangle (2.6,1.9);
    \node[below right] at (-2.6,1.9) {$E$};
    \begin{scope}
      \clip (-2.6,-1.9) rectangle (2.6,1.9);
      \fill[omThm!20] (-2.6,-1.9) rectangle (2.6,1.9);
      \begin{scope}
        \clip (-0.7,0) circle (1.1);
        \fill[white] (0.7,0) circle (1.1);
      \end{scope}
    \end{scope}
    \draw[omThm, thick] (-0.7,0) circle (1.1);
    \draw[omThm, thick] (0.7,0) circle (1.1);
    \node[omThm] at (-1.35,0.5) {$A$};
    \node[omThm] at (1.35,0.5) {$B$};
  \end{scope}
\end{tikzpicture}

{\small Hukum De Morgan dalam gambar: daerah yang diarsir di sebelah
kiri adalah $\overline{A \cup B} = \overline A \cap \overline B$
(segala sesuatu di luar kedua cakram); di sebelah kanan,
$\overline{A \cap B} = \overline A \cup \overline B$ (segala sesuatu
kecuali daerah tumpang-tindih yang berbentuk lensa). Sebuah diagram
bukanlah bukti, tetapi ia membuat bukti pengejaran unsur pada
\cref{prop:b1:logic:setalgebra} mustahil salah diingat.}
\end{omfigure}

\section{Pemetaan}

\begin{definition}[Pemetaan, peta, prapeta]\label{def:b1:logic:map}
Sebuah \emph{pemetaan}\index{pemetaan} (atau \emph{fungsi})
$f \colon E \to F$ mengaitkan setiap unsur $x$ pada \omterm{def:b1:logic:sets}{himpunan} $E$
(\emph{daerah asal}) dengan tepat satu unsur $f(x)$ pada \omterm{def:b1:logic:sets}{himpunan} $F$
(\emph{daerah kawan}). Untuk $A
\subseteq E$ dan $B \subseteq F$:
\[
f(A) = \{f(x) : x \in A\} \subseteq F,
\qquad
f^{-1}(B) = \{x \in E : f(x) \in B\} \subseteq E
\]
berturut-turut adalah \emph{peta langsung} dari $A$ dan \emph{prapeta}
\index{prapeta} dari $B$. \emph{Komposisi} dari $f \colon E \to F$
dan $g \colon F \to G$ adalah $g \circ f \colon E \to G$, $x \mapsto
g(f(x))$.
\end{definition}

\begin{remark}
Notasi $f^{-1}(B)$ \emph{tidak} mengandaikan adanya \omterm{def:b1:logic:map}{pemetaan} invers:
$f^{-1}(B)$ terdefinisi untuk setiap $f$. \omterm{def:b1:logic:map}{Prapeta} berperilaku lebih
baik daripada peta: $f^{-1}$ mengawetkan gabungan, irisan dan
komplemen, sedangkan $f(A \cap A') \subseteq f(A) \cap f(A')$ dapat
berupa inklusi sejati (\cref{exo:b1:logic:8}).
\end{remark}

\begin{example}[Menghitung peta dan prapeta]\label{ex:b1:logic:images}
Misalkan $f \colon \R \to \R$, $x \mapsto x^2$. Maka:
\[
f\bigl(\intcc{-1}{2}\bigr) = \intcc04, \qquad
f^{-1}\bigl(\intcc14\bigr) = \intcc{-2}{-1} \cup \intcc12, \qquad
f^{-1}(\{-1\}) = \emptyset .
\]
Untuk yang pertama: setiap $x \in \intcc{-1}2$ memenuhi
$x^2 \in \intcc04$, dan setiap $y \in \intcc04$ tercapai sebagai
$y = (\sqrt y)^2$ dengan $\sqrt y
\in \intcc02 \subseteq \intcc{-1}2$ --- perhatikan bahwa petanya
\emph{bukan} $\intcc14 = \{(-1)^2, 2^2\}$: peta sebuah selang tidak
dihitung dari titik ujungnya saja. Untuk yang kedua:
$1 \leq x^2 \leq 4 \iff
1 \leq \abs x \leq 2$, yang terpecah menjadi dua potong. Yang ketiga
memperlihatkan bahwa \omterm{def:b1:logic:map}{prapeta} boleh kosong --- $f^{-1}(B)$ selalu
bermakna, sekecil apa pun irisan $B$ dengan peta itu. Akhirnya
perhatikan pada contoh ini gejala inklusi sejati dari catatan di atas:
dengan $A = \intcc{-1}0$ dan $A' = \intcc01$, kita peroleh
$f(A \cap A') = f(\{0\}) = \{0\}$, sedangkan
$f(A) \cap f(A') = \intcc01$.
\end{example}

\begin{definition}[Injektif, surjektif, bijektif]\label{def:b1:logic:inj}
Sebuah \omterm{def:b1:logic:map}{pemetaan} $f \colon E \to F$ disebut:
\begin{itemize}
  \item \emph{injektif}\index{injektif} bila unsur yang berbeda
        mempunyai peta yang berbeda: $\forall x, x' \in E,\ f(x) = f(x') \implies
        x = x'$;
  \item \emph{surjektif}\index{surjektif} bila setiap unsur $F$
        tercapai: $\forall y \in F,\ \exists x \in E,\ f(x) = y$;
  \item \emph{bijektif}\index{bijektif} bila keduanya berlaku, yakni
        setiap $y \in F$ mempunyai tepat satu \omterm{def:b1:logic:map}{prapeta}.
\end{itemize}
\end{definition}

\begin{theorem}[Pemetaan invers]\label{thm:b1:logic:inverse}
\omterm{def:b1:logic:map}{Pemetaan} $f \colon E \to F$ \omterm{def:b1:logic:inj}{bijektif} jika dan hanya jika ada \omterm{def:b1:logic:map}{pemetaan}
$g \colon F \to E$ dengan $g \circ f = \mathrm{id}_E$ dan $f \circ g =
\mathrm{id}_F$. Dalam hal itu $g$ tunggal; ia ditulis $f^{-1}$ dan
disebut \emph{invers} dari $f$, dan $f^{-1}$ sendiri \omterm{def:b1:logic:inj}{bijektif} dengan
$(f^{-1})^{-1} = f$.
\end{theorem}

\begin{proof}
($\Rightarrow$) Jika $f$ \omterm{def:b1:logic:inj}{bijektif}, setiap $y \in F$ mempunyai tepat
satu \omterm{def:b1:logic:map}{prapeta}; definisikan $g(y)$ sebagai \omterm{def:b1:logic:map}{prapeta} itu. Maka
$f(g(y)) = y$ menurut konstruksinya, dan $g(f(x)) = x$ karena $x$
adalah \emph{satu-satunya} \omterm{def:b1:logic:map}{prapeta} $f(x)$.

($\Leftarrow$) Andaikan $g$ semacam itu ada. Jika $f(x) = f(x')$,
penerapan $g$ memberikan $x = x'$: jadi $f$ \omterm{def:b1:logic:inj}{injektif}. Untuk
$y \in F$, unsur $x = g(y)$
memenuhi $f(x) = y$: jadi $f$ \omterm{def:b1:logic:inj}{surjektif}.

Ketunggalan: jika $g$ dan $h$ sama-sama memenuhi, maka
$g = g \circ \mathrm{id}_F
= g \circ (f \circ h) = (g \circ f) \circ h = h$. Akhirnya pasangan
kesamaan itu simetris dalam $f$ dan $g$, sehingga $g = f^{-1}$
\omterm{def:b1:logic:inj}{bijektif} dengan invers $f$.
\end{proof}

\begin{example}[Menghitung invers dalam praktik]\label{ex:b1:logic:inversecomp}
Misalkan $f \colon \R \to \intoo0{+\infty}$, $f(x) = \eu^{2x+1}$.
Untuk membalikkannya, selesaikan $y = f(x)$ bagi $y > 0$ yang
diberikan:
\[
y = \eu^{2x+1} \iff \ln y = 2x + 1 \iff x = \frac{\ln y - 1}2 ,
\]
setiap langkahnya dapat dibalik pada daerah yang telah diumumkan.
Perhitungan itu sekaligus menyerahkan segalanya: untuk setiap $y$ pada
daerah kawan ada tepat satu penyelesaian $x$, sehingga $f$ \omterm{def:b1:logic:inj}{bijektif},
dan
\[
f^{-1} \colon \intoo0{+\infty} \to \R,
\qquad
f^{-1}(y) = \frac{\ln y - 1}2 .
\]
Pemeriksaan cepat atas kedua komposisinya ($f^{-1}(f(x)) = \frac{(2x+1) -
1}2 = x$ dan $f(f^{-1}(y)) = \eu^{\ln y} = y$) membenarkan kriteria
\cref{thm:b1:logic:inverse}. Inti gagasannya: ``selesaikan untuk $x$
dan awasi ekuivalensinya'' sekaligus merupakan bukti keberadaan, bukti
ketunggalan, dan rumusnya --- tetapi itu hanya berhasil bila daerah
kawannya diumumkan dengan benar ($f$ \emph{tidak}
\omterm{def:b1:logic:inj}{surjektif} ke $\R$).
\end{example}

\begin{proposition}[Komposisi dan ketiga sifat itu]\label{prop:b1:logic:comp}
Misalkan $f \colon E \to F$ dan $g \colon F \to G$.
\begin{enumerate}
  \item Jika $f$ dan $g$ \omterm{def:b1:logic:inj}{injektif} (masing-masing \omterm{def:b1:logic:inj}{surjektif}, \omterm{def:b1:logic:inj}{bijektif}),
        maka demikian pula $g \circ f$; dan dalam kasus \omterm{def:b1:logic:inj}{bijektif}
        berlaku $(g \circ f)^{-1} = f^{-1} \circ
        g^{-1}$.
  \item Jika $g \circ f$ \omterm{def:b1:logic:inj}{injektif}, maka $f$ \omterm{def:b1:logic:inj}{injektif}. Jika
        $g \circ f$ \omterm{def:b1:logic:inj}{surjektif}, maka $g$ \omterm{def:b1:logic:inj}{surjektif}.
\end{enumerate}
\end{proposition}

\begin{proof}
(1) Jika $g(f(x)) = g(f(x'))$, keinjektifan $g$ memberikan
$f(x) = f(x')$, lalu keinjektifan $f$ memberikan $x = x'$. Jika
$z \in G$, kesurjektifan $g$ memberikan $y$ dengan $g(y) = z$, lalu
kesurjektifan $f$ memberikan $x$ dengan
$f(x) = y$, sehingga $g(f(x)) = z$. Dalam kasus \omterm{def:b1:logic:inj}{bijektif} kita periksa
langsung bahwa $f^{-1} \circ g^{-1}$ adalah invers dua sisi dari
$g \circ f$, dan ketunggalan pada \cref{thm:b1:logic:inverse} menutup
buktinya.

(2) Jika $f(x) = f(x')$ maka $g(f(x)) = g(f(x'))$, dan keinjektifan
$g \circ f$ memberikan $x = x'$. Jika $z \in G$, kesurjektifan
$g \circ f$ memberikan $x$ dengan $g(f(x)) = z$: lalu $y = f(x)$
memenuhi $g(y) = z$.
\end{proof}

\begin{example}[Butir (2) sudah tajam]\label{ex:b1:logic:compsharp}
Pada \cref{prop:b1:logic:comp}~(2), kesimpulannya tidak dapat
diperkuat: $g \circ f$ \omterm{def:b1:logic:inj}{bijektif} \emph{tidak} memaksa $f$
\omterm{def:b1:logic:inj}{surjektif} atau $g$ \omterm{def:b1:logic:inj}{injektif}. Ambil $E = G = \{1\}$, $F = \{1,
2\}$, dengan $f(1) = 1$ dan $g(1) = g(2) = 1$: maka $g \circ f =
\mathrm{id}_E$ \omterm{def:b1:logic:inj}{bijektif}, namun $f$ melewatkan unsur $2$ dan
$g$ meleburkan kedua unsurnya. Moralnya adalah aturan pembukuan yang
tepat: informasi komposisi mengalir ke \omterm{def:b1:logic:map}{pemetaan} \emph{dalam} untuk
keinjektifan dan ke \omterm{def:b1:logic:map}{pemetaan} \emph{luar} untuk kesurjektifan, tidak
pernah sebaliknya. (\cref{exo:b1:logic:9} membangun gejala yang sama
dengan \omterm{def:b1:logic:sets}{himpunan} tak hingga, tempat gejala itu menjadi mesin di balik
invers sepihak.)
\end{example}

\begin{example}\label{ex:b1:logic:injexamples}
$f \colon \R \to \R$, $x \mapsto x^2$ tidak \omterm{def:b1:logic:inj}{injektif} ($f(-1) =
f(1)$) dan tidak pula \omterm{def:b1:logic:inj}{surjektif} ($-1$ tak punya \omterm{def:b1:logic:map}{prapeta}). Dengan
membatasi daerah asal dan daerah kawan, $f \colon \R_+ \to \R_+$,
$x \mapsto x^2$ \omterm{def:b1:logic:inj}{bijektif}, dengan invers $y \mapsto \sqrt y$.
Keinjektifan atau kesurjektifan sebuah \omterm{def:b1:logic:map}{pemetaan} bergantung pada daerah
asal dan daerah kawan yang diumumkan, bukan hanya pada rumusnya.
\end{example}

\section{Relasi}

\begin{definition}[Relasi ekuivalensi]\label{def:b1:logic:equiv}
Sebuah \emph{relasi biner} $\mathcal{R}$ pada \omterm{def:b1:logic:sets}{himpunan} $E$ disebut
\emph{relasi ekuivalensi}\index{relasi ekuivalensi} bila ia bersifat
\emph{refleksif} ($x \mathbin{\mathcal{R}} x$ untuk setiap $x$),
\emph{simetris} ($x \mathbin{\mathcal{R}} y \implies y
\mathbin{\mathcal{R}} x$) dan \emph{transitif} ($x
\mathbin{\mathcal{R}} y$ dan $y \mathbin{\mathcal{R}} z$
mengakibatkan $x
\mathbin{\mathcal{R}} z$). \emph{Kelas ekuivalensi} dari $x$ adalah
$\mathrm{cl}(x) = \{y \in E : x \mathbin{\mathcal{R}} y\}$.
\end{definition}

\begin{example}[Memeriksa ketiga aksioma]\label{ex:b1:logic:fracpart}
Pada $\R$, tetapkan $x \mathbin{\mathcal{R}} y$ bila $x - y \in \Z$.
\emph{Refleksif:} $x - x = 0 \in \Z$. \emph{Simetris:} jika $x - y
\in \Z$ maka $y - x = -(x - y) \in \Z$. \emph{Transitif:} jika $x -
y \in \Z$ dan $y - z \in \Z$, maka $x - z = (x - y) + (y - z) \in
\Z$ (jumlah dua bilangan bulat). Jadi $\mathcal R$ adalah \omterm{def:b1:logic:equiv}{relasi
ekuivalensi}, dan $\mathrm{cl}(x) = x + \Z = \{x + k : k \in \Z\}$:
setiap kelas memuat tepat satu wakil di $\intco01$, yaitu
\emph{bagian pecahan}-nya. Sebaliknya, relasi ``$\abs{x - y}
\leq 1$'' pada $\R$ bersifat refleksif dan simetris tetapi
\emph{tidak} transitif ($0 \mathbin{\mathcal R} 1$ dan
$1 \mathbin{\mathcal R}
2$, padahal $\abs{0 - 2} > 1$): kedekatan tidak menular, dan tidak ada
\omterm{thm:b1:logic:partition}{partisi} ke dalam kelas --- contoh penyangkal yang berguna untuk
diingat ketika pemeriksaan aksioma mulai terasa rutin.
\end{example}

\begin{theorem}[Kelas-kelasnya membentuk partisi]\label{thm:b1:logic:partition}
Misalkan $\mathcal{R}$ \omterm{def:b1:logic:equiv}{relasi ekuivalensi} pada $E$. Maka kelas-kelas
ekuivalensinya tak kosong, saling lepas atau sama, dan gabungannya
adalah $E$: kelas-kelas itu membentuk \emph{partisi}\index{partisi}
dari $E$. Sebaliknya, setiap \omterm{thm:b1:logic:partition}{partisi} dari $E$ muncul dengan cara ini
dari tepat satu \omterm{def:b1:logic:equiv}{relasi ekuivalensi} (``berada pada potongan yang sama'').
\end{theorem}

\begin{proof}
$x \in \mathrm{cl}(x)$ menurut kerefleksifan, jadi setiap kelas tak
kosong dan gabungannya $E$. Andaikan
$\mathrm{cl}(x) \cap \mathrm{cl}(y) \neq \emptyset$, katakanlah $z$
terletak di keduanya. Maka $x \mathbin{\mathcal{R}} z$ dan $y
\mathbin{\mathcal{R}} z$, sehingga menurut kesimetrisan dan
ketransitifan $x
\mathbin{\mathcal{R}} y$. Selanjutnya untuk sebarang
$t \in \mathrm{cl}(y)$, ketransitifan memberikan
$t \in \mathrm{cl}(x)$, dan begitu pula sebaliknya: kedua
kelas itu sama. Untuk arah sebaliknya, misalkan $(E_i)_{i \in I}$
sebuah \omterm{thm:b1:logic:partition}{partisi} dari $E$ dan definisikan $x \mathbin{\mathcal S} y$
sebagai ``ada potongan yang memuat $x$ dan $y$ sekaligus''.
\emph{Refleksif:} $x$ terletak pada suatu potongan, yang lalu memuat
$x$ dua kali. \emph{Simetris:} syarat pendefinisinya simetris dalam
$x$ dan $y$. \emph{Transitif:} jika $x, y \in E_i$ dan
$y, z \in E_j$, maka $y \in E_i \cap E_j$, sehingga $E_i = E_j$
(potongan yang berbeda saling lepas) dan $x, z$ berbagi satu potongan.
Kelas-$\mathcal S$ dari $x$ tepat sama dengan potongan yang memuat
$x$, jadi kelas-kelasnya adalah potongan yang diberikan itu. Akhirnya
relasinya ditentukan oleh kelas-kelasnya: dua \omterm{def:b1:logic:equiv}{relasi ekuivalensi}
dengan kelas yang sama merelasikan pasangan yang sama, karena
masing-masing merelasikan $x$ dan $y$ tepat ketika $y$ termasuk dalam
kelas $x$ --- dari sanalah klaim ketunggalannya.
\end{proof}

\begin{example}\label{ex:b1:logic:congruence}
Pada $\Z$, kekongruenan modulo $n$ ($x \equiv y \pmod n$ bila $n$
membagi $x - y$) adalah \omterm{def:b1:logic:equiv}{relasi ekuivalensi}; kelas-kelasnya adalah $n$
\omterm{def:b1:logic:sets}{himpunan} bilangan bulat yang bersisa tertentu bila dibagi $n$. Contoh
ini menjelma menjadi ring $\Z/n\Z$ pada \cref{ch:b1:structures}.
\end{example}

\begin{definition}[Relasi urutan]\label{def:b1:logic:order}
Relasi $\preceq$ pada $E$ disebut \emph{urutan}\index{relasi urutan}
bila ia refleksif, \emph{antisimetris} ($x \preceq y$ dan $y
\preceq x$ mengakibatkan $x = y$) serta transitif. Urutan itu
\emph{total} bila setiap dua unsur dapat dibandingkan, dan
\emph{parsial} bila tidak. Unsur $M \in A \subseteq E$ disebut
\emph{unsur terbesar} dari $A$ bila $a \preceq M$ untuk setiap
$a \in A$; unsur terbesar (dan terkecil) bersifat tunggal bila ada.
\end{definition}

\begin{example}\label{ex:b1:logic:orders}
$(\R, \leq)$ terurut total. $(\mathcal{P}(E), \subseteq)$ terurut
parsial begitu $E$ mempunyai dua unsur: $\{a\}$ dan
$\{b\}$ tak dapat dibandingkan. \omterm{def:b1:logic:sets}{Himpunan} bagian $A = \{\{a\}, \{b\}\}$
dari $\mathcal{P}(\{a,b\})$ tak mempunyai unsur terbesar, namun
mempunyai batas atas $\{a, b\}$: pembedaan antara unsur terbesar dan
batas atas muncul kembali, untuk $\R$, pada \cref{ch:b1:reals}.
\end{example}

\begin{example}[Dua urutan pada kisi $\N^2$]\label{ex:b1:logic:lexorder}
Pada pasangan bilangan cacah, bandingkan komponen demi komponen:
$(a, b) \preceq (a',
b')$ bila $a \leq a'$ \emph{dan} $b \leq b'$ (\emph{\omterm{def:b1:logic:order}{urutan} hasil
kali}). Ini memang \omterm{def:b1:logic:order}{urutan} --- setiap aksiomanya diwarisi koordinat
demi koordinat --- tetapi \omterm{def:b1:logic:order}{urutan} parsial: $(1, 3)$ dan $(2, 0)$ tak
terbandingkan. Kini bandingkan seperti kamus: $(a, b)
\preceq_{\mathrm{lex}} (a', b')$ bila $a < a'$, atau $a = a'$ dan $b
\leq b'$ (\emph{\omterm{def:b1:logic:order}{urutan} leksikografis}). Ketransitifannya menuntut
pemeriksaan dua kasus tetapi tetap berlaku, dan kini setiap dua
pasangan dapat dibandingkan: \omterm{def:b1:logic:order}{urutannya} total. Kedua \omterm{def:b1:logic:order}{urutan} itu
memeringkat \omterm{def:b1:logic:sets}{himpunan} yang sama secara berbeda
--- $(0, 100) \preceq_{\mathrm{lex}} (1, 0)$ padahal \omterm{def:b1:logic:order}{urutan} hasil
kali tak berkata apa-apa --- sebuah pengingat bahwa \omterm{def:b1:logic:order}{urutan} adalah
struktur yang kita \emph{pilih}, bukan sifat dari \omterm{def:b1:logic:sets}{himpunannya}.
Perbandingan leksikografis juga merupakan kiat baku untuk melebur
beberapa kriteria pengurutan menjadi satu.
\end{example}

\begin{remark}[Selingan: ukuran sebagai bijeksi]\label{rem:b1:logic:interlude}
Sebuah tema yang mengalir diam-diam sepanjang bab ini pantas disorot:
bijeksi adalah gagasan matematikawan tentang ``ukuran yang sama''.
Untuk \omterm{def:b1:logic:sets}{himpunan} hingga hal ini menjelma menjadi kalkulus pencacahan
pada \cref{ch:b1:counting}, tempat setiap rumus diam-diam merupakan
sebuah bijeksi; untuk \omterm{def:b1:logic:sets}{himpunan} tak hingga hal itu menjelma menjadi
soal akhir pekan di bawah, tempat $\N$, $\Q$ dan $\R$ ternyata
mempunyai ukuran yang sungguh-sungguh berbeda. Kamus yang sama muncul
dua kali lagi dalam jilid ini dalam bentuk yang lebih halus: barisan
(\cref{ch:b1:seq}) tidak lain adalah \omterm{def:b1:logic:map}{pemetaan} $\N \to \R$, sehingga
\omterm{def:b1:logic:statement}{pernyataan} tentang barisan adalah \omterm{def:b1:logic:statement}{pernyataan} tentang sebuah \omterm{def:b1:logic:sets}{himpunan}
\omterm{def:b1:logic:map}{pemetaan}; dan aljabar linear akan mengukur ruang vektor bukan dengan
bijeksi melainkan dengan bijeksi \emph{linear}, yang keberadaannya
dikendalikan oleh satu bilangan tunggal, yaitu dimensi
(\cref{ch:b1:findim}). Setiap kali sebuah ``kesamaan'' yang baru
muncul --- ekuipotensi, isomorfisma grup
(\cref{ch:b1:structures}), isomorfisma linear --- pola
\cref{thm:b1:logic:inverse} berulang: kesamaan adalah \omterm{def:b1:logic:map}{pemetaan} yang
punya invers dan menghormati struktur.
\end{remark}

\begin{remark}[Di mana bab ini dipakai]\label{rem:b1:logic:whereused}
Di mana-mana --- tetapi beberapa tempat pantas ditandai. Senam tiga
kuantor pada \cref{ex:b1:logic:limit} adalah santapan sehari-hari
\cref{ch:b1:seq,ch:b1:continuity}: setiap bukti limit adalah permainan
melawan $\varepsilon$ yang sembarang. \omterm{def:b1:logic:equiv}{Kelas ekuivalensi} muncul kembali
sebagai kelas kekongruenan $\Z/n\Z$ pada \cref{ch:b1:structures},
tempat \omterm{thm:b1:logic:partition}{partisi} \cref{thm:b1:logic:partition} memperoleh struktur
aljabar tersendiri. \omterm{def:b1:logic:order}{Relasi urutan}, batas atas dan batas atas terkecil
menjadi jantung aksiomatis $\R$ pada \cref{ch:b1:reals}.
Injeksi, surjeksi dan bijeksi kembali sebagai \omterm{def:b1:logic:map}{pemetaan} linear pada
\cref{ch:b1:linmaps}, tempat keinjektifan dapat diuji pada satu
vektor saja (kernelnya); dan soal akhir pekan di bawah mengubah
gagasan bijeksi yang telanjang menjadi teori tentang \emph{ukuran
\omterm{def:b1:logic:sets}{himpunan} tak hingga}, yang kesimpulannya (keterbilangan $\Q$,
ketakterbilangan $\R$) muncul lagi pada
\cref{ch:b1:reals,ch:b1:topology}.
\end{remark}

\section{Latihan}

\begin{exercise}[$\star$]\label{exo:b1:logic:1}
Tuliskan negasi setiap \omterm{def:b1:logic:statement}{pernyataan} berikut, tanpa memakai kata ``bukan'':
\begin{enumerate}
  \item $\forall x \in \R,\ \exists y \in \R,\ x + y > 0$;
  \item $\exists x \in \R,\ \forall y \in \R,\ xy = 0$;
  \item $\forall \varepsilon > 0,\ \exists \delta > 0,\ \forall x \in
        \R,\ \abs{x} \leq \delta \implies \abs{f(x)} \leq \varepsilon$
        (untuk sebuah \omterm{def:b1:logic:map}{pemetaan} tetap $f \colon \R \to \R$).
\end{enumerate}
Kemudian tentukan apakah \omterm{def:b1:logic:statement}{pernyataan} (1) dan (2) benar.
\end{exercise}

\begin{exercise}[$\star$]\label{exo:b1:logic:2}
Misalkan $P, Q$ \omterm{def:b1:logic:statement}{pernyataan}. Dengan tabel kebenaran, buktikan bahwa
$\lnot(P \implies Q) \iff P \land (\lnot Q)$, lalu simpulkan negasi
dari: ``jika sebuah fungsi dapat diturunkan maka ia kontinu''.
\end{exercise}

\begin{exercise}[$\star$]\label{exo:b1:logic:3}
Buktikan dengan kontraposisi: untuk $x \in \R$, jika $x^3 + x \geq 2$
maka $x \geq 1$. Kemudian buktikan dengan kontradiksi: tidak ada
bilangan real positif tegas yang terkecil.
\end{exercise}

\begin{exercise}[$\star$]\label{exo:b1:logic:4}
Buktikan dengan induksi bahwa untuk setiap $n \in \N$:
\begin{enumerate}
  \item $\sum_{k=0}^{n} 2^k = 2^{n+1} - 1$;
  \item $4^n + 5$ habis dibagi $3$.
\end{enumerate}
\end{exercise}

\begin{exercise}[$\star$]\label{exo:b1:logic:5}
Temukan cacat pada ``bukti'' berikut bahwa semua pensil berwarna sama.
\emph{Misalkan $P(n)$: ``pada setiap \omterm{def:b1:logic:sets}{himpunan} berisi $n$ pensil, semua
pensil berwarna sama''. $P(1)$ jelas. Andaikan $P(n)$ lalu ambil $n+1$
pensil; dengan membuang yang terakhir, $n$ pensil pertama sewarna;
dengan membuang yang pertama, $n$ pensil terakhir sewarna; jadi semua
$n+1$ pensil sewarna.}
\end{exercise}

\begin{exercise}[$\star$]\label{exo:b1:logic:6}
Misalkan $A, B, C$ \omterm{def:b1:logic:sets}{himpunan} bagian dari $E$. Buktikan:
\begin{enumerate}
  \item $A \setminus B = A \cap \overline{B}$;
  \item $(A \cup B) \setminus C = (A \setminus C) \cup (B \setminus
        C)$;
  \item $A \subseteq B \iff A \cup B = B \iff A \cap B = A$.
\end{enumerate}
\end{exercise}

\begin{exercise}[$\star\star$]\label{exo:b1:logic:7}
Untuk setiap \omterm{def:b1:logic:map}{pemetaan} berikut, tentukan (dengan bukti) apakah ia
\omterm{def:b1:logic:inj}{injektif}, \omterm{def:b1:logic:inj}{surjektif}, atau \omterm{def:b1:logic:inj}{bijektif}:
\begin{enumerate}
  \item $f \colon \N \to \N$, $n \mapsto n + 1$;
  \item $g \colon \Z \to \Z$, $n \mapsto n + 1$;
  \item $h \colon \R \setminus \{1\} \to \R$, $x \mapsto
        \frac{x+1}{x-1}$.
\end{enumerate}
Untuk $h$, sesuaikan daerah kawannya agar ia \omterm{def:b1:logic:inj}{bijektif}, lalu hitung
inversnya.
\end{exercise}

\begin{exercise}[$\star\star$]\label{exo:b1:logic:8}
Misalkan $f \colon E \to F$, dan misalkan $A, A' \subseteq E$ serta
$B, B' \subseteq
F$.
\begin{enumerate}
  \item Buktikan $f^{-1}(B \cap B') = f^{-1}(B) \cap f^{-1}(B')$ dan
        $f(A \cup A') = f(A) \cup f(A')$.
  \item Buktikan $f(A \cap A') \subseteq f(A) \cap f(A')$ lalu berikan
        contoh yang inklusinya sejati.
  \item Buktikan: $f$ \omterm{def:b1:logic:inj}{injektif} jika dan hanya jika $f(A \cap A') = f(A)
        \cap f(A')$ untuk setiap $A, A'$.
\end{enumerate}
\end{exercise}

\begin{exercise}[$\star\star$]\label{exo:b1:logic:9}
Misalkan $f \colon E \to F$ dan $g \colon F \to E$ memenuhi $g \circ f =
\mathrm{id}_E$. Buktikan bahwa $f$ \omterm{def:b1:logic:inj}{injektif} dan $g$ \omterm{def:b1:logic:inj}{surjektif}.
Berikan contoh yang $f$ maupun $g$-nya tidak \omterm{def:b1:logic:inj}{bijektif}.
\end{exercise}

\begin{exercise}[$\star\star$]\label{exo:b1:logic:10}
Pada $\R$, definisikan $x \mathbin{\mathcal{R}} y \iff x^2 - y^2 = x - y$.
Buktikan bahwa $\mathcal{R}$ adalah \omterm{def:b1:logic:equiv}{relasi ekuivalensi} lalu jelaskan
kelas ekuivalensi setiap bilangan real $x$. Kelas mana yang
beranggotakan tepat satu unsur?
\end{exercise}

\begin{exercise}[$\star\star\star$]\label{exo:b1:logic:11}
(Cantor) Misalkan $E$ sebuah \omterm{def:b1:logic:sets}{himpunan}. Buktikan bahwa tidak ada
surjeksi dari $E$ pada $\mathcal{P}(E)$. \emph{Petunjuk: diberikan
$f \colon E \to
\mathcal{P}(E)$, tinjau $D = \{x \in E : x \notin f(x)\}$.}
\end{exercise}

\begin{exercise}[$\star\star\star$]\label{exo:b1:logic:12}
Misalkan $f \colon E \to F$ sebuah \omterm{def:b1:logic:map}{pemetaan}. Definisikan $\Phi \colon \mathcal{P}(F)
\to \mathcal{P}(E)$ dengan $\Phi(B) = f^{-1}(B)$.
\begin{enumerate}
  \item Buktikan bahwa $f$ \omterm{def:b1:logic:inj}{surjektif} jika dan hanya jika $\Phi$
        \omterm{def:b1:logic:inj}{injektif}.
  \item Buktikan bahwa $f$ \omterm{def:b1:logic:inj}{injektif} jika dan hanya jika $\Phi$
        \omterm{def:b1:logic:inj}{surjektif}.
\end{enumerate}
\end{exercise}

\section{Soal: Membandingkan ketakhinggaan}

\begin{problem}\label{pb:b1:logic:1}
Kapan dua \omterm{def:b1:logic:sets}{himpunan} dikatakan mempunyai ``banyak unsur yang sama''?
Jawaban Cantor --- ketika ada bijeksi di antara keduanya --- ternyata
dapat dipakai bahkan untuk \omterm{def:b1:logic:sets}{himpunan} tak hingga, dan jawaban itu
memecah ketakhinggaan menjadi ukuran yang sungguh berbeda. Soal ini
membangun seluruh perkakasnya dari definisi telanjang bab ini: teorema
Cantor--Schröder--Bernstein (dua injeksi memproduksi sebuah bijeksi),
keterbilangan $\Q$, ketakterbilangan $\R$ lewat argumen diagonal, dan
kesimpulan Cantor tahun 1874 yang mencengangkan: \emph{\omterm{pb:b1:logic:1}{bilangan
transenden} itu ada, dan sangat banyak}, tanpa menunjukkan satu pun.
\index{himpunan terbilang}\index{teorema Cantor--Schroder--Bernstein}
Di sepanjang soal ini, untuk \omterm{def:b1:logic:sets}{himpunan} $E$ dan $F$, tulis $E \preceq F$
bila ada injeksi dari $E$ ke dalam $F$, dan $E \approx F$ (``$E$ dan
$F$ \emph{ekuipoten}'') bila ada bijeksi dari $E$
pada $F$.

\medskip
\noindent\textbf{Bagian I --- Kosakata perbandingan.}
\begin{enumerate}
  \item Tunjukkan bahwa $\approx$ berperilaku seperti \omterm{def:b1:logic:equiv}{relasi
        ekuivalensi}: $E \approx E$; jika $E \approx F$ maka
        $F \approx E$; jika $E \approx F$ dan $F \approx G$ maka
        $E \approx G$. (Kutip secara persis
        \cref{thm:b1:logic:inverse} dan
        \cref{prop:b1:logic:comp}.)
  \item Tunjukkan bahwa $\preceq$ transitif, dan bahwa injeksi
        $f \colon E \to F$ selalu melahirkan $E \approx f(E)$.
  \item Misalkan $E \neq \emptyset$. Tunjukkan bahwa $E \preceq F$
        jika dan hanya jika ada surjeksi dari $F$ pada $E$.
  \item Periksa bahwa $n \mapsto n + 1$ adalah bijeksi dari $\N$ pada
        $\N^* = \N \setminus \{0\}$, dan bahwa
        \[
        \sigma(n) = \frac n2 \ \ (n \text{ genap}), \qquad
        \sigma(n) = -\frac{n+1}2 \ \ (n \text{ ganjil})
        \]
        adalah bijeksi dari $\N$ pada $\Z$. Jadi membuang satu titik,
        atau melipatgandakan ke arah negatif, tidak mengubah ukuran
        $\N$.
\end{enumerate}

\medskip
\noindent\textbf{Bagian II --- Teorema
Cantor--Schröder--Bernstein.} Misalkan $f \colon E \to F$ dan
$g \colon F \to E$ dua injeksi. Definisikan
\[
C_0 = E \setminus g(F), \qquad C_{n+1} = g\bigl(f(C_n)\bigr)
\ \ (n \in \N), \qquad C = \bigcup_{n \in \N} C_n,
\]
lalu misalkan $h \colon E \to F$ memetakan $x \in C$ ke $f(x)$, dan
$x \notin C$ ke satu-satunya $y \in F$ dengan $g(y) = x$.
\begin{enumerate}[resume]
  \item Periksa bahwa $h$ terdefinisi dengan baik: jika $x \notin C$
        maka $x \in
        g(F)$, dan unsur $y$ dengan $g(y) = x$ bersifat tunggal.
  \item Tunjukkan bahwa $g\bigl(f(C)\bigr) = \bigcup_{n \geq 1} C_n
        \subseteq C$. (Peta langsung berkomutasi dengan gabungan:
        \cref{exo:b1:logic:8}.)
  \item Tunjukkan bahwa $h$ \omterm{def:b1:logic:inj}{injektif}. (Tiga kasus; pada kasus campuran
        $x \in C$, $x' \notin C$, tunjukkan bahwa $h(x) = h(x')$ akan
        memaksa $x' \in g(f(C)) \subseteq C$.)
  \item Tunjukkan bahwa $h$ \omterm{def:b1:logic:inj}{surjektif}: diberikan $y \in F$, bedakan
        kasus $g(y) \notin C$ dan $g(y) \in C_n$ untuk suatu $n
        \geq 1$ (mengapa $g(y) \in C_0$ mustahil?), lalu tunjukkan
        sebuah \omterm{def:b1:logic:map}{prapeta} dari $y$ pada masing-masing kasus.
  \item Simpulkan dengan \emph{teorema
        Cantor--Schröder--Bernstein}: jika $E \preceq F$ dan
        $F \preceq E$, maka $E
        \approx F$. Berikan satu kalimat komentar tentang apa yang
        membuat \omterm{def:b1:logic:statement}{pernyataan} ini tidak sepele.
  \item Dua penerapan. (a) Tunjukkan $\intcc01 \approx \intoo01$.
        (b) Tunjukkan bahwa $\varphi(p, q) = 2^p(2q + 1) - 1$
        mendefinisikan bijeksi dari $\N \times \N$ pada $\N$ ---
        keinjektifan lewat argumen paritas, kesurjektifan lewat
        induksi kuat (\cref{thm:b1:logic:induction}). Jadi
        $\N \times \N \approx
        \N$: bidang titik bulat tidak lebih besar daripada
        garis.
\end{enumerate}

\medskip
\noindent\textbf{Bagian III --- \omterm{pb:b1:logic:1}{Himpunan terbilang}.} Sebuah \omterm{def:b1:logic:sets}{himpunan}
$E$ disebut \emph{paling banyak terbilang} bila $E \preceq \N$, dan
\emph{terbilang} bila $E \approx \N$.
\begin{enumerate}[resume]
  \item Tunjukkan bahwa setiap \omterm{def:b1:logic:sets}{himpunan} bagian tak hingga
        $A \subseteq \N$ bersifat terbilang. (Definisikan $\varphi(n)$
        secara rekursif sebagai unsur terkecil dari
        $A \setminus \{\varphi(0), \dots,
        \varphi(n-1)\}$; tunjukkan bahwa $\varphi$ naik tegas,
        memenuhi $\varphi(n) \geq n$, dan mencapai setiap unsur
        $A$.)
  \item Simpulkan bahwa sebuah \omterm{def:b1:logic:sets}{himpunan} paling banyak terbilang jika
        dan hanya jika ia hingga atau terbilang, lalu amati bahwa
        pertanyaan 9 memberi jalan pintas: jika $E \preceq \N$ dan
        $\N \preceq E$, maka $E$ terbilang.
  \item Tunjukkan bahwa jika $E$ dan $F$ paling banyak terbilang, maka
        demikian pula $E
        \times F$. Simpulkan bahwa $\Z \times \N^*$ terbilang.
  \item Tunjukkan bahwa $\Q$ terbilang. (Injeksikan $\Q$ ke dalam
        $\Z \times
        \N^*$ dengan menuliskan setiap bilangan rasional dalam bentuk
        paling sederhana berpenyebut positif --- ketunggalan penyajian
        itu dibuktikan pada \cref{ch:b1:arith}; lalu terapkan
        pertanyaan 12.)
  \item Tunjukkan bahwa gabungan terbilang dari \omterm{def:b1:logic:sets}{himpunan} yang paling
        banyak terbilang tetap paling banyak terbilang: jika setiap
        $E_n$ ($n \in \N$) paling banyak terbilang, maka demikian pula
        $\bigcup_{n \in \N} E_n$. (Kirim $x$ ke pasangan
        $(n, f_n(x))$ dengan $n$ indeks \emph{terkecil} yang memenuhi
        $x \in E_n$.)
  \item Tunjukkan bahwa \omterm{def:b1:logic:sets}{himpunan} semua \omterm{def:b1:logic:sets}{himpunan} bagian \emph{hingga}
        dari $\N$ bersifat terbilang. (Petakan \omterm{def:b1:logic:sets}{himpunan} bagian hingga
        $F$ ke $\sum_{i \in F} 2^i$;
        buktikan keinjektifannya dengan membandingkan unsur terbesar
        tempat kedua \omterm{def:b1:logic:sets}{himpunan} hingga itu berbeda, memakai
        $\sum_{k=0}^{m-1} 2^k = 2^m - 1$
        dari \cref{exo:b1:logic:4}.)
\end{enumerate}

\medskip
\noindent\textbf{Bagian IV --- Diagonalisasi.} Misalkan
$\{0,1\}^{\N}$ menyatakan \omterm{def:b1:logic:sets}{himpunan} semua \omterm{def:b1:logic:map}{pemetaan}
$u \colon \N \to \{0, 1\}$, yakni \omterm{def:b1:logic:sets}{himpunan} barisan biner.
\begin{enumerate}[resume]
  \item Konstruksikan bijeksi antara $\mathcal{P}(\N)$ dan
        $\{0,1\}^{\N}$ (fungsi indikator).
  \item (Argumen diagonal) Misalkan $\Phi \colon \N \to
        \{0,1\}^{\N}$ sebuah \omterm{def:b1:logic:map}{pemetaan} sembarang. Tinjau barisan $d$
        yang didefinisikan oleh $d(n) = 1 - \Phi(n)(n)$. Tunjukkan
        bahwa $d$ tidak berada pada peta $\Phi$, lalu simpulkan bahwa
        $\{0,1\}^{\N}$ \emph{tidak} paling banyak terbilang. Jelaskan
        dalam satu kalimat mengapa, lewat pertanyaan 17, hal ini
        tepat merupakan teorema Cantor
        (\cref{exo:b1:logic:11}) untuk $E = \N$.
  \item Terima --- sebagaimana sudah dikenal sejak sekolah, dan
        ditegakkan secara ketat pada \cref{ch:b1:reals} --- bahwa
        setiap $x \in
        \intco01$ mempunyai satu ekspansi desimal \emph{sejati} yang
        tunggal $x =
        0.d_1 d_2 d_3\dots$ (yaitu yang tidak berakhir dengan untaian
        tak hingga angka $9$). Diberikan sebarang barisan
        $(x_n)_{n \geq 1}$ unsur
        $\intco01$, konstruksikan $x \in \intco01$ dengan
        $x \neq x_n$ untuk setiap $n$: pilih angka ke-$n$-nya sama
        dengan $5$ bila angka ke-$n$
        dari $x_n$ berbeda dari $5$, dan $6$ bila tidak. Berikan
        alasan cermat bahwa $x$ sejati dan menghindari setiap $x_n$,
        lalu simpulkan bahwa $\intco01$ tidak paling banyak terbilang.
  \item Simpulkan bahwa $\R$ tak terbilang, dan bahwa \omterm{def:b1:logic:sets}{himpunan} $\R
        \setminus \Q$ berisi bilangan irasional juga tak terbilang.
        Dalam pengertian apa persisnya ``kebanyakan'' bilangan real
        bersifat irasional?
\end{enumerate}

\medskip
\noindent\textbf{Bagian V --- Teorema Cantor tahun 1874: \omterm{pb:b1:logic:1}{bilangan
transenden} itu ada.} Bilangan real $x$ disebut \emph{aljabar}
\index{bilangan aljabar} bila $P(x) = 0$ untuk suatu polinomial tak
nol $P$ berkoefisien bulat, dan \emph{transenden}
\index{bilangan transenden} bila tidak. Terima untuk bagian ini ---
hal ini dibuktikan pada \cref{ch:b1:poly} --- bahwa polinomial tak nol
berderajat $n$ mempunyai paling banyak $n$ akar real.
\begin{enumerate}[resume]
  \item Tunjukkan bahwa setiap bilangan rasional bersifat aljabar, dan
        carilah polinomial berkoefisien bulat yang secara eksplisit
        menolkan $\sqrt 2$ dan $\sqrt 2 + \sqrt 3$.
  \item Untuk $n \in \N$ yang tetap, tunjukkan bahwa \omterm{def:b1:logic:sets}{himpunan}
        polinomial berderajat paling tinggi $n$ dengan koefisien bulat
        bersifat terbilang. (Injeksikan ia ke dalam $\Z^{n+1}$ lalu
        berinduksi pada $n$ dengan pertanyaan 13.)
  \item Simpulkan bahwa \omterm{def:b1:logic:sets}{himpunan} \emph{semua} polinomial berkoefisien
        bulat bersifat terbilang.
  \item Buktikan \emph{teorema Cantor tentang \omterm{pb:b1:logic:1}{bilangan aljabar}}:
        \omterm{def:b1:logic:sets}{himpunan} $\mathcal{A}$ berisi bilangan real aljabar bersifat
        terbilang.
  \item Simpulkan: bilangan real transenden itu ada, dan \omterm{def:b1:logic:sets}{himpunan}
        \omterm{pb:b1:logic:1}{bilangan transenden} tak terbilang. Lalu tariklah kesimpulan
        atas seluruh soal ini dalam beberapa kalimat: rantai
        $\N \approx
        \Z \approx \Q \approx \mathcal{A}$, lompatan tegas ke $\R
        \approx$ (pada dasarnya) $\mathcal{P}(\N)$, tempat setiap
        perkakas (Cantor--Schröder--Bernstein, gabungan terbilang,
        argumen diagonal) menjadi penentu --- dan pukulan filosofis
        dari pembuktian bahwa ada tak terbilang banyaknya \omterm{pb:b1:logic:1}{bilangan
        transenden} tanpa menyebut satu pun. (Membuktikan sebuah
        bilangan \emph{tertentu} seperti $\pi$ transenden jauh lebih
        sulit dan berada di luar jilid ini.)

\end{enumerate}
\end{problem}
